\chapter{Design Philosophy and House Style}
\label{ch:design-philosophy}

Every worked example in the rest of this book silently assumes a set of conventions that CNA
enforces on itself through its \texttt{CLAUDE.md} porting rules and \texttt{CHECKLIST.md}.
This chapter makes those conventions explicit, once, so that later chapters can simply use
them.

\section{The prime directive: match FNA, not taste}

CNA's porting rules name a single authoritative behavioral reference: a local FNA source
tree. The instruction is blunt: ``Do not treat old CNA code or AI-generated stubs as
authoritative if they conflict with the FNA reference API.'' Concretely, this means:

\begin{itemize}
  \item Class names, struct names, enum names, method names, operator names, and constant
        names must match XNA/FNA \emph{exactly}. A method is not renamed to be ``more
        C\texttt{++}-like'' if that would diverge from the reference.
  \item Behavior is matched over personal preference: packed value layouts (\cnaclass{Color}
        stores its channels as \texttt{AABBGGRR}), clamping behavior, integer cast behavior,
        operator behavior, method overloads, default values, and exception behavior where
        practical, all follow FNA even when a different choice would look cleaner in
        isolation.
  \item C\texttt{++} member declaration order mirrors the original C\# source order where
        practical, specifically to keep diff-based review tractable against the reference.
\end{itemize}

This is the single idea that makes the rest of the house style legible: every convention
below exists to answer the same question --- \emph{how do you represent a C\# concept in
C\texttt{++} without losing the ability to diff against the original?}

\section{Two namespaces, two jobs}

CNA code lives in exactly one of two namespace families, and the rule for which one is
absolute:

\begin{center}
\begin{tabular}{ll}
\toprule
\textbf{Namespace} & \textbf{Contains} \\
\midrule
\cnans{Microsoft::Xna::Framework::*} & Real XNA 4.0 types, matching the reference exactly \\
\cnans{CNA::*} & Project-specific extensions, helpers, internal backends \\
\bottomrule
\end{tabular}
\end{center}

Original XNA types are never moved into the \texttt{CNA} namespace, and the reverse is
equally firm: anything that is not part of the real XNA~4.0 API must not silently live inside
\cnans{Microsoft::Xna::Framework} unmarked. That is what the \texttt{NOXNA} marker exists for.

\section{The \texttt{NOXNA} tag: a compiler-enforced honesty check}
\label{sec:noxna}

\texttt{NOXNA} is an empty marker macro, defined in \texttt{include/CNA/CNAHelper.hpp}, whose
only job is to visibly tag a declaration inside the XNA namespace tree as \emph{not} part of
the original XNA~4.0 API --- a rumble-and-gyro extension on \cnaclass{GamePad}, a
\texttt{GetTypeName()} override, a raw joystick API, a haptics subsystem. Two concrete
examples from the real codebase:

\begin{lstlisting}
// A real XNA 4.0 override, unmarked:
[[nodiscard]] bool Equals(const Color& other) const;

// A CNA-only addition living in the same class, explicitly marked:
NOXNA [[nodiscard]] const std::string& GetTypeName() const override;
\end{lstlisting}

What makes this more than a documentation convention is a compiler-enforced check described in
Chapter~\ref{ch:what-is-cna}, precisely: \texttt{CNAHelper.hpp}'s own definition of
\texttt{NOXNA} is conditional on a second, distinct macro, \texttt{CNA\_STRICT\_XNA\_API} --- in
a normal build \texttt{NOXNA} expands to nothing at all (it is documentation only); only when
\texttt{CNA\_STRICT\_XNA\_API} is defined does it expand to
\texttt{[[deprecated("NOXNA: not part of the XNA 4.0 API surface")]]}. That define is never set
for the normal \texttt{CNA} library target --- it is scoped to exactly two small standalone
CMake targets in \texttt{cmake/Harnesses.cmake}, each compiled with
\texttt{-Werror=deprecated-declarations} on top of it, so a stray \texttt{NOXNA}-tagged call
anywhere those targets' translation units can see turns into a hard build failure rather than a
warning that could be scrolled past.

\begin{sourcenote}
\textbf{A worked example: the check that verifies the check.} A positive-only test --- ``the
real API still compiles clean'' --- cannot by itself prove the mechanism would actually catch a
real leak; it is equally consistent with a check that always passes. CNA's own
\texttt{tools/\allowbreak{}devices/\allowbreak{}Strict\allowbreak{}Xna\allowbreak{}Api\allowbreak{}Surface\allowbreak{}Check.cpp} (compiled as the
\texttt{cna\allowbreak{}\_strict\allowbreak{}\_xna\allowbreak{}\_api\allowbreak{}\_check} target,
run as the \texttt{StrictXnaApiSurfaceCheck\allowbreak{}\_Compile\allowbreak{}\_Run} CTest) is
that positive direction, and its own header comment names every \texttt{NOXNA}-tagged
\cnaclass{Accelerometer} / \cnaclass{Sensors} member it deliberately avoids calling. The negative
direction --- proving a real leak genuinely fails --- is a second, separate file,
\texttt{tools/\allowbreak{}devices/\allowbreak{}Strict\allowbreak{}Xna\allowbreak{}Api\allowbreak{}Surface\allowbreak{}Leak\allowbreak{}Check.cpp}, whose entire body is one deliberate
violation:
\begin{lstlisting}
Microsoft::Devices::Sensors::Accelerometer accelerometer;
accelerometer.InjectSyntheticSensorUpdate(1.0f, 0.0f, 0.0f); // NOXNA -- must fail to compile here.
\end{lstlisting}
This compiles into an \texttt{EXCLUDE\_FROM\_ALL} target
(\texttt{cna\allowbreak{}\_strict\allowbreak{}\_xna\allowbreak{}\_api\allowbreak{}\_leak\allowbreak{}\_check}),
so it never participates in an ordinary build. It is invoked by exactly one CTest:
\texttt{Strict\allowbreak{}Xna\allowbreak{}Api\allowbreak{}Surface\allowbreak{}Leak\allowbreak{}Check\allowbreak{}\_Must\allowbreak{}Fail\allowbreak{}To\allowbreak{}Compile}, whose
\texttt{WILL\_FAIL TRUE} property inverts the usual pass/fail meaning --- the test
\emph{passes} only if the build command it wraps \emph{fails}. If a future change to the
\texttt{NOXNA} macro or to \texttt{CNA\_STRICT\_XNA\_API}'s wiring ever let this
deliberately-broken file compile, this specific test would flip from passing to failing ---
catching a regression in the safety net itself, not merely in the API surface the net is meant
to protect.
\end{sourcenote}

\section{C\# properties, translated consistently}

C\# auto-properties (\texttt{public byte R \{ get; set; \}}) become a fixed pair of C\texttt{++}
methods, never a public field, unless the surrounding type has already established a
field-style convention:

\begin{lstlisting}
// C#:  public byte R { get; set; }

// C++:
[[nodiscard]] bytecs getRProperty() const;
void setRProperty(bytecs value);
\end{lstlisting}

\section{SharpRuntime type aliases: never a raw C\texttt{++} primitive in the XNA surface}

Wherever C\# source uses a .NET primitive type name, CNA code uses the corresponding
\texttt{sharp-runtime} alias rather than a raw C\texttt{++} fundamental type, specifically so
that the XNA-facing surface keeps a visible, greppable link back to its .NET origin:

\begin{center}
\begin{tabular}{llll}
\toprule
\textbf{C\# type} & \textbf{sharp-runtime alias} & \textbf{Underlying C\texttt{++} type} \\
\midrule
\texttt{byte}   & \texttt{bytecs} / \texttt{Byte}   & \texttt{uint8\_t} \\
\texttt{sbyte}  & \texttt{sbytecs} / \texttt{SByte} & \texttt{int8\_t} \\
\texttt{short}  & \texttt{shortcs} / \texttt{Int16} & \texttt{int16\_t} \\
\texttt{ushort} & \texttt{ushortcs} / \texttt{UInt16} & \texttt{uint16\_t} \\
\texttt{int}    & \texttt{intcs} / \texttt{Int32}   & \texttt{int32\_t} \\
\texttt{uint}   & \texttt{uintcs} / \texttt{UInt32} & \texttt{uint32\_t} \\
\texttt{long}   & \texttt{longcs} / \texttt{Int64}  & \texttt{int64\_t} \\
\texttt{ulong}  & \texttt{ulongcs} / \texttt{UInt64} & \texttt{uint64\_t} \\
\texttt{float}  & \texttt{Single}                   & \texttt{float} \\
\texttt{string} & \texttt{String}                   & \texttt{std::string} \\
\texttt{char}   & \texttt{charcs}                   & \texttt{char16\_t} \\
\bottomrule
\end{tabular}
\end{center}

If a needed alias does not yet exist, the rule is to add a minimal stub to
\texttt{sharp-runtime} first --- never to reach for a raw C\texttt{++} type directly in the
XNA API surface as a shortcut. This is the same missing-dependency discipline covered in
\S\ref{sec:missing-deps} below.

\section{Events, as \texttt{System::EventHandler<T>}}

C\# events and delegates are modeled uniformly through \texttt{sharp-runtime}'s
\texttt{System::EventHandler<T>}, never through an ad-hoc project-specific callback type:

\begin{lstlisting}
// Declaration, matching C# "public event EventHandler<T> Name;"
System::EventHandler<ExitingEventArgs> Exiting;

// Subscription
Exiting += [](System::Object* sender, const ExitingEventArgs& e) { /* ... */ };

// Raising
Exiting.Raise(this, args);   // or Exiting.Invoke(this, args);
\end{lstlisting}

\section{Interfaces as abstract base classes, IDisposable as a pattern}

C\# interface relationships become C\texttt{++} abstract base classes --- for example
\cnaclass{Color} implementing \texttt{IEquatable<Color>} and \texttt{IPackedVector} in C\#
becomes \texttt{struct Color : public Graphics::PackedVector::IPackedVectorT<UInt32>} in
CNA. Where an exact mapping is not practical, the rule is to implement equivalent behavior
and document the intentional deviation in the pull request description --- explicitly
\emph{not} as a source comment, keeping the deviation visible to reviewers without cluttering
the header for every future reader.

\texttt{IDisposable} follows a fixed pattern backed by \texttt{sharp-runtime}'s
\texttt{System::IDisposable}:

\begin{lstlisting}
class Foo : public System::IDisposable {
public:
    void Dispose() override;
protected:
    void Dispose(bool disposing);   // only when the pattern requires it
};
\end{lstlisting}

with an \texttt{isDisposed\_} guard checked before any operation, throwing
\texttt{std::runtime\_error} on use-after-dispose.

\section{Visibility is chosen, not defaulted}

C\# \texttt{internal} members do not automatically become C\texttt{++} \texttt{public} just
because C\texttt{++} lacks an \texttt{internal} keyword --- they become \texttt{private},
\texttt{protected}, live in a detail/internal namespace, or are omitted from the port
entirely. The house style calls out a concrete anti-example: a C\# \texttt{internal
DebugDisplayString} helper should not become a public C\texttt{++} API method just because it
was easy to leave public.

\section{File structure mirrors namespace, always}

Declarations live in \texttt{.hpp} under \texttt{include/}; implementation lives in
\texttt{.cpp} under \texttt{src/}; and the directory path mirrors the namespace path exactly:

\begin{lstlisting}
include/Microsoft/Xna/Framework/Graphics/Texture2D.hpp
src/Microsoft/Xna/Framework/Graphics/Texture2D.cpp
\end{lstlisting}

Non-template implementations do not live in headers. The same mirroring rule is why this
book can describe ``the Graphics namespace'' and ``the 106 headers under
\texttt{include/\allowbreak{}Microsoft/\allowbreak{}Xna/\allowbreak{}Framework/\allowbreak{}Graphics/}''
as the same thing throughout Part~III.

\section{No backward-compatibility shims}

If correcting an API breaks an old demo, the demo is fixed --- CNA does not carry a
convenience alias just to keep outdated call sites compiling. The house style gives its own
canonical wrong/right pair:

\begin{lstlisting}
// Wrong: an old-CNA-era shortcut kept around for compatibility
inline const Color CornflowerBlue(/* ... */);

// Correct: the real, single source of truth
Color::CornflowerBlue
\end{lstlisting}

\section{Documentation: full Doxygen, or nothing}

Every public method, constructor, property accessor, operator, and constant in every
\texttt{.hpp} file must carry a Doxygen block --- bare \texttt{///} line comments are never
acceptable on a public declaration (they remain acceptable only for brief inline notes inside
a method body). The intent of the original C\# XML doc comments (\texttt{<summary>},
\texttt{<param>}, \texttt{<returns>}) is preserved, often verbatim from FNA where the wording
already fits, but never marked with a comment like ``taken from FNA'' --- provenance belongs
in the porting record, not scattered through the header.

\section{Missing dependencies: stub correctly, never invent}
\label{sec:missing-deps}

When a file being ported references a type that does not exist yet, the rule is layered:
if the missing type belongs to the .NET runtime (\texttt{System.*} or a primitive alias), it
is added to \texttt{sharp-runtime} first, following the same minimal-stub discipline
described in Chapter~\ref{ch:ecosystem-map}. Otherwise, a minimal, \emph{correctly named}
stub is added directly in CNA, in the correct final namespace, sufficient only to compile ---
never a large unrelated system built just to satisfy one missing dependency. Every stub
introduced this way is reported in the task or PR description, so it is never silently
forgotten as ``real'' once the build goes green.

\section{The per-file porting checklist, in miniature}

\texttt{CHECKLIST.md} governs the full requirements for porting one FNA source file to CNA;
the house rules single out its non-negotiable minimum:

\begin{itemize}
  \item An SPDX \texttt{MS-PL} license identifier at the top of both the \texttt{.hpp} and
        the \texttt{.cpp}.
  \item \texttt{\#include "CNA/CNAHelper.hpp"} in the header, whenever \texttt{NOXNA} is used
        anywhere in that file.
  \item Every method body checked line-by-line against the FNA equivalent, with every
        intentional deviation documented in a source comment (not silently absorbed).
  \item Every concrete class inheriting \texttt{System::Object} overrides
        \texttt{GetTypeName()}, tagged \texttt{NOXNA}, returning the fully-qualified .NET
        name (e.g.\ \texttt{"Microsoft.Xna.Framework.Game"}).
  \item A unit test for every public method, operator, and constant --- with out-ref
        overloads tested separately from their value-returning counterparts, and static
        factory methods (\cnaclass{CreateFromPoints}, \cnaclass{CreateMerged},
        \cnaclass{CreateFromSphere}, \ldots) each getting a dedicated test.
\end{itemize}

A file is ``done in one pass'' --- the house style explicitly rejects a
``make-and-forget-partially'' workflow where checklist items are deferred to a later,
unscheduled pass.

\section{Why this chapter earns its place in a ``bible''}

None of the conventions above are unique or clever in isolation --- type-alias tables and
Doxygen mandates exist in a thousand C\texttt{++} style guides. What is unusual is how tightly
they are all in service of one measurable goal: keeping a large, multi-backend,
multi-platform C\texttt{++} codebase auditable against a real external reference (FNA, real
XNA~4.0, \texttt{xna4-spec}) at every layer --- namespace, visibility, member order, type
alias, and a compiler flag that turns ``did we leak a non-XNA type into the XNA surface''
into a build error instead of a hoped-for code-review catch. The rest of this book takes
these conventions as read; when a later chapter says a class ``matches FNA exactly,'' this is
the mechanism that claim rests on.
